\abstract{РЕФЕРАТ}

Объем работы равен \formbytotal{lastpage}{страниц}{е}{ам}{ам}. Работа содержит \formbytotal{figurecnt}{иллюстраци}{ю}{и}{й}, \formbytotal{tablecnt}{таблиц}{у}{ы}{}, \arabic{bibcount} библиографических источников и \formbytotal{числоПлакатов}{лист}{}{а}{ов} графического материала. Количество приложений – 2. Графический материал представлен в приложении А. Фрагменты исходного кода представлены в приложении Б.

Перечень ключевых слов: система управления, бизнес, ресторан, касса, программно-информационная система, база данных, заказ, отчёт, автоматизация, оптимизация, классы, сотрудники, прибыль, технологические карты, повар, администратор, официант, управляющий.

Объектом разработки является программно-информационная система управления ресторанном бизнесом, оптимизирующая рабочих процессов заведений общественного питания.

Целью выпускной квалификационной работы является разработка программы, направленной на оптимизацию ресторанного бизнеса.

В процессе создания программы были выделены основные объекты путем создания информационных блоков; использованы классы и методы модулей, обеспечивающие работу с сущностями предметной области; а также корректную работу программно-информационной системы, разработаны разделы, содержащие информацию о функциональности программы в различных процессах и кассовых отчётах.

Программа была написала на языке Python с использованием системы-управления базами данных PostgreSQL.

\selectlanguage{english}
\abstract{ABSTRACT}
  
The volume of work is \formbytotal{lastpage}{page}{}{s}{s}. The work contains \formbytotal{figurecnt}{illustration}{}{s}{s}, \formbytotal{tablecnt}{table}{}{s}{s}, \arabic{bibcount} bibliographic sources and \formbytotal{числоПлакатов}{sheet}{}{s}{s} of graphic material. The number of applications is 2. The graphic material is presented in annex A. The layout of the site, including the connection of components, is presented in annex B.

The list of keywords: management system, business, restaurant, cash register, software and information system, database, order, report, automation, optimization, classes, employees, reception, technological maps, cook, administrator, waiter, manager.

The object of the development is a software and information management system for the restaurant business that optimizes the work processes of catering establishments.

The purpose of the final qualification is to develop a program aimed at optimizing the restaurant business.

In the process of creating the program, the main entities were identified by creating information blocks, classes and module methods were used to work with the entities of the subject area, as well as the correct operation of the software information system, sections were developed containing information about the functionality of the program in various processes and cash reports.

The program was written in Python using a PostgreSQL database management system.
\selectlanguage{russian}
